\usepackage[left=30mm, right=10mm, top=20mm, bottom=20mm]{geometry}
% XeLaTeX: Unicode і системні шрифти (збирати через xelatex або latexmk)
\usepackage{fontspec}
\usepackage[ukrainian]{babel}
\usepackage{setspace}
\usepackage{graphicx}
\usepackage{subcaption}
\usepackage{titlesec}
\usepackage{pdflscape}
\usepackage{amsmath}
\usepackage{enumitem}
\usepackage{icomma}
\usepackage{tikz}
\usepackage{tikz-cd}      % діаграми залежностей модулів (tikzcd)
\usetikzlibrary{babel}    % щоб `"` у підписах стрілок не конфліктував з babel (ukrainian)
\usepackage{xcolor}
\usepackage{float}
\usepackage{tcolorbox}
\usepackage{listings}
\usepackage{url}         % links in the bibliography break across lines
\usepackage[nounderscore]{syntax}      % BNF: середовище grammar, <нетермінал>, \lit{термінал}, \alt (syntax-mdw)
\tcbuselibrary{listings, skins, breakable}

% #### Шрифти ####
% Основний шрифт, схожий на Computer Modern; для Times: \setmainfont{TeX Gyre Termes}
\setmainfont{CMU Serif}[Scale=0.94]
% Iosevka Xast лежить у fonts/ (з лігатурою .\ у вигляді лямбди)
\newfontfamily\xastfont[
    Path=fonts/, Extension=.ttf,
    UprightFont=IosevkaXast-Medium, ItalicFont=IosevkaXast-Medium-Italic,
    BoldFont=IosevkaXast-Medium, BoldItalicFont=IosevkaXast-Medium-Italic,
]{IosevkaXast}
\setmonofont[
    Path=fonts/, Extension=.ttf,
    UprightFont=IosevkaXast-Medium, ItalicFont=IosevkaXast-Medium-Italic,
    BoldFont=IosevkaXast-Medium, BoldItalicFont=IosevkaXast-Medium-Italic,
    Scale=MatchLowercase,
]{IosevkaXast}

% #### Лістинги ####
% Лістинги коду: мова Xast, колір, лігатури Iosevka Xast.
%
%   \xast|fn id = .\x -> x|        рядковий код (це \lstinline: межі будь-який символ, якого нема в коді;
%                                    у фігурних дужках код не може містити `{` і `}`; не вкладати в аргументи макросів)
%   \code{ім'я}                    рядковий моноширинний текст
%   \begin{xastcode} ... \end{xastcode}              блок коду
%   \begin{xastlisting}{Назва}  ... \end{xastlisting}  нумерований лістинг (Лістинг 1.1)
%   \xastinput{файл.xst}  /  \xastinput[title=...]{файл.xst}   код із файлу
%   ключі: numbered (номери рядків), title={...}, label=lst:назва
% Те саме є для Haskell (haskellcode) і C (ccode).

\definecolor{codeKeyword}{HTML}{1F4FBF}
\definecolor{codeType}{HTML}{7B3F98}
\definecolor{codeAnnotation}{HTML}{B35900}
\definecolor{codeComment}{HTML}{6E6E6E}
\definecolor{codeString}{HTML}{0A7F7A}
\definecolor{codeNumber}{HTML}{B0306A}

% Лігатури Iosevka: listings друкує символи окремо, тому оператори
% підставляються цілим блоком, і шрифт будує з них лігатуру.
\newcommand{\codeligature}[1]{\mbox{\xastfont #1}}

\lstdefinelanguage{Xast}{
    morekeywords={module, exports, use, as, type, fn, extern, system, match, with,
        let, in, if, then, else, copy, boxed, number, concatenative},
    morekeywords=[2]{Int, Long, Short, Byte, UInt, ULong, UShort, UByte, Size, USize,
        Float, Double, Bool, Char, String, List, Vec, Option, Result},
    morekeywords=[3]{@Copyable, @Pod, @Drop, @Singleton, @Intrinsic, @Inline,
        @SharedVariant, @Parallel, @Deprecated},
    alsoletter={@'},
    sensitive=true,
    morecomment=[l]{--},
    morestring=[b]",
    literate=
        {.\\}{{\codeligature{.\textbackslash}}}2
        {->}{{\codeligature{->}}}2
        {<-}{{\codeligature{<-}}}2
        {=>}{{\codeligature{=>}}}2
        {<>}{{\codeligature{<>}}}2
        {|>}{{\codeligature{|>}}}2
        {==}{{\codeligature{==}}}2
        {!=}{{\codeligature{!=}}}2
        {<=}{{\codeligature{<=}}}2
        {>=}{{\codeligature{>=}}}2
        {::}{{\codeligature{::}}}2
        {..}{{\codeligature{..}}}2
}

\lstdefinestyle{xastcodestyle}{
    basicstyle=\xastfont\small,
    keywordstyle=\color{codeKeyword},
    keywordstyle=[2]\color{codeType},
    keywordstyle=[3]\color{codeAnnotation},
    commentstyle=\color{codeComment}\itshape,
    stringstyle=\color{codeString},
    numberstyle=\tiny\color{codeComment},
    columns=fixed,
    keepspaces=true,
    showstringspaces=false,
    extendedchars=false,
    texcl=true,   % коментарі друкує LaTeX: так працює кирилиця
    tabsize=4,
    breaklines=true,
    postbreak=\mbox{\textcolor{codeComment}{$\hookrightarrow$}\space},
}

\lstset{style=xastcodestyle}

% Загальні налаштування рамки
\tcbset{
    codebox/.style={
        listing only,
        breakable,
        enhanced,
        colback=black!4,
        colframe=black!60,
        boxrule=0.4pt,
        arc=1pt,
        left=4pt, right=4pt, top=3pt, bottom=3pt,
        fonttitle=\normalfont\small,
        coltitle=black,
        colbacktitle=black!12,
        before skip=8pt, after skip=8pt,
    },
    numbered/.style={
        listing options={numbers=left, numbersep=8pt, numberstyle=\scriptsize\color{codeComment}, xleftmargin=8pt},
    },
}

\newtcblisting{xastcode}[1][]{codebox, listing options={language=Xast}, #1}
\newtcblisting{haskellcode}[1][]{codebox, listing options={language=Haskell}, #1}
\newtcblisting{ccode}[1][]{codebox, listing options={language=C}, #1}

% Нумерований лістинг (не плаває, розривається між сторінками):
% \begin{xastlisting}[label=lst:id]{Назва} ... \end{xastlisting}
\newtcblisting[auto counter, number within=chapter]{xastlisting}[2][]{
    codebox, listing options={language=Xast},
    title={Лістинг~\thetcbcounter~--- #2},
    title after break={Лістинг~\thetcbcounter{} (продовження)},
    #1
}

% Код із файлу
\newcommand{\xastinput}[2][]{%
    \tcbinputlisting{codebox, listing options={language=Xast}, listing file={#2}, #1}%
}

% Рядковий код
\newcommand{\xast}{\lstinline[language=Xast]}
\newcommand{\code}[1]{\texttt{#1}}


\titleformat{\chapter}
    {\normalfont\Large\bfseries\centering}
    {\thechapter}{1em}{}

\titleformat{\section}
    {\normalfont\large\bfseries}
    {\thesection}{1em}{}

\titleformat{\subsection}
    {\normalfont\large\bfseries}
    {\thesubsection}{1em}{}

% #### Списки: без зайвих проміжків, ромб як маркер першого рівня ####
\setlist{itemsep=0pt, topsep=0pt}
\setlist[itemize,1]{label=$\diamond$}

\onehalfspacing
\setlength{\parindent}{0cm}
\setlength{\parskip}{0.5em}
